\section{Yann LeCun y el marco JEPA}

Al hablar de Yann LeCun, la descripción de Xie Saining es muy vívida: el LeCun de internet es un «luchador», mientras que el LeCun de la vida real es «una persona muy, muy buena». Este contraste no es raro en sí mismo, pero resulta especialmente interesante en el caso de LeCun, porque él siempre ha desempeñado simultáneamente el papel de polemista dentro de la comunidad científica y el de estabilizador dentro de su organización. Por un lado, defiende con firmeza en público su juicio sobre las limitaciones de los LLM, sin cambiar de postura fácilmente por conveniencia de la narrativa corporativa; por otro, tiene un interés personal y un sentido de la vida muy fuertes: construye aeromodelos, hace astrofotografía, escucha música electrónica y jazz, navega en velero. Estas aficiones lo hacen parecer alguien que vive en una adolescencia prolongada de manera permanente. Aquella frase de He Kaiming, \textit{``LeCun es una persona cuya adolescencia de 16 años se ha extendido hasta los 65 años''}, no es una broma, sino una síntesis precisa de una curiosidad y un espíritu lúdico desbordantes.

\begin{knowledgebox}{La estructura de personalidad de LeCun}
El LeCun que ve Xie Saining no es la imagen única del «ganador del Premio Turing», sino la de una persona que sigue construyendo cosas con las manos, que sigue siendo curiosa y que sigue debatiendo en público de manera constante. Lo que tienen en común aficiones como los aeromodelos, la astronomía, el jazz y el velero es que todas exigen una dedicación de largo plazo, en las que coexisten el sentido técnico y el sentido estético, y todas requieren que la persona respete a la vez las restricciones de la realidad y conserve la emoción de explorar lo desconocido.
\end{knowledgebox}

Esta estructura de personalidad es coherente con su estilo de gestión. Xie Saining dice que LeCun dirige a su equipo como quien navega en velero: la mayor parte del tiempo hay que confiar plenamente en que cada persona actúe bien desde su propio puesto; pero en cuanto se percibe una desviación en el viento, en la corriente o en la actitud del casco, hay que corregir cuanto antes, en lugar de esperar a que el problema crezca para tirar de vuelta con fuerza. Para una organización de investigación, esta es una filosofía de gestión muy madura, porque a las personas verdaderamente creativas no se les puede supervisar de manera excesiva, pero si se les deja hacer por completo sin control, el equipo tiende a dispersarse en intereses parciales. El núcleo de este método de gestión al estilo del velero consiste en que la organización mantenga el sentido de dirección sin, por ello, asfixiar de manera prematura la capacidad de cada persona para descubrir problemas por sí misma.

JEPA se entiende precisamente en este contexto. Xie Saining subraya de manera especial que JEPA no es un modelo, sino un océano muy amplio. No es «una combinación de módulos de un artículo en particular», sino un marco general sobre la representación, la predicción y la planificación. Por eso resume su propio proceso de comprensión de JEPA en tres etapas: cuestionar JEPA, comprender JEPA, convertirse en JEPA. Esta forma de decirlo suena a broma, pero capta con precisión la ruta cognitiva de muchos marcos de vanguardia. Al principio uno tiende a pensar que es demasiado abstracto, demasiado amplio, como un concepto en el que casi todo cabe; pero cuando se ven una y otra vez las charlas de LeCun y se coloca el marco de nuevo en el contexto de los modelos del mundo, uno se da cuenta poco a poco de que la amplitud de JEPA no es vaguedad, sino que en realidad intenta cubrir toda la cadena que va de la comprensión del mundo (\textit{world understanding}) a la predicción (\textit{prediction}) y de ahí a la planificación (\textit{planning}).

\begin{importantbox}{Las tres etapas de comprensión de JEPA}
\begin{tabularx}{\textwidth}{>{\raggedright\arraybackslash}p{0.2\textwidth} Y}
\toprule
Etapa & Estado típico \\
\midrule
Cuestionar JEPA & Se percibe como demasiado abstracto, no como una receta «ejecutable de inmediato», difícil de alinear con un benchmark existente. \\
Comprender JEPA & Se empieza a percibir que describe un objetivo de aprendizaje de más alto nivel: no ajustarse a píxeles o tokens, sino aprender representaciones abstractas con valor predictivo sobre el estado del mundo. \\
Convertirse en JEPA & Deja de tratarse como un método aislado y pasa a pensarse la representación, la predicción, la planificación y la acción dentro de un mismo marco unificado. \\
\bottomrule
\end{tabularx}
\end{importantbox}

Este pensamiento por marcos también sustenta la crítica de larga data de LeCun a los LLM. Xie Saining cita una frase muy teatral: \textit{``los LLM acabarán por marchitarse... los viejos soldados no mueren, solo se marchitan''}. Esto no significa que los modelos de lenguaje vayan a fallar de repente, sino que no llegarán a ser el fundamento definitivo para construir sistemas de inteligencia general. La razón es clara: si un sistema aprende sobre todo a partir de secuencias de texto, queda en gran medida confinado a «descripciones del mundo ya comprimidas por otros», sin haber construido de manera suficiente su propio modelo del mundo frente a la realidad. La ambición de JEPA consiste precisamente en sortear esta limitación y aprender representaciones más cercanas a la estructura misma del mundo.

\begin{warningbox}{Tomar al LLM como destino final es una ilusión propia de una etapa}
LeCun y Xie Saining no niegan el enorme valor de los LLM, pero ambos se oponen a extrapolar directamente la capacidad actual de los modelos de lenguaje como si fuera el destino final. El lenguaje es una herramienta de abstracción muy poderosa, pero no es la totalidad de la realidad. Si un sistema depende sobre todo del lenguaje para alinearse con el mundo, sigue avanzando apoyado en muletas. Las muletas pueden ayudarte a ponerte de pie, pero no necesariamente bastan para que corras dentro del mundo real.
\end{warningbox}

LeCun resulta convincente para Xie Saining también porque siempre conserva una integridad propia de científico. Una observación muy importante de la entrevista es esta: aunque la empresa a la que pertenece, la tendencia del momento y las preferencias del mercado de capitales empujen todos hacia una narrativa de «acercarse a los LLM», LeCun no modifica en público su juicio fundamental sobre la ruta tecnológica solo porque la organización lo necesite. Esta firmeza no siempre resulta cómoda, pero es sumamente importante para una comunidad de investigación, porque si hasta los juicios más centrales pueden reescribirse a voluntad según la moda del momento, entonces la llamada dirección de investigación termina degradándose en una simple prolongación del discurso de mercadotecnia.

El hecho de que Xie Saining haya visto una y otra vez las charlas de LeCun también muestra que JEPA no es el tipo de marco que «se entiende en cuanto se escucha». Exige que quien investiga vaya construyendo poco a poco un hábito de abstracción de más alto nivel: no apresurarse a preguntar a qué ranking, a qué función de pérdida o a qué módulo corresponde, sino preguntar primero qué objetivo de aprendizaje está describiendo en realidad. Precisamente porque este proceso de comprensión es lento, JEPA se presta a que en las primeras etapas se le juzgue como algo vago; y precisamente porque es lo bastante amplio, es posible que en el futuro llegue a albergar un sistema más completo de comprensión y planificación del mundo.

Xie Saining menciona además que ha visto las charlas de LeCun entre diez y veinte veces, y que cada vez obtiene algo nuevo. Esta frase explica muy bien el umbral de comprensión de JEPA. Un marco verdaderamente amplio no se asimila por completo la primera vez que se escucha; requiere que uno vuelva a él una y otra vez con la experiencia acumulada en cada etapa. En cierto sentido, esto es muy coherente con su manera de entender la investigación: una buena teoría no se limita a darte una conclusión, sino que sigue ofreciendo nueva visibilidad conforme cambia tu propia capacidad.

Por último, hay una comparación muy cálida. Xie Saining dice: \textit{``Yann es una batería enorme, él me empuja, y yo espero transmitir esa energía hacia adelante''}. Esta frase muestra que LeCun no es para él solo un mentor intelectual, sino también una fuente de energía. Un líder académico verdaderamente fuerte no solo produce ideas, sino que también eleva de manera conjunta el potencial de quienes lo rodean, para que estén dispuestos a seguir apostando por problemas cada vez más difíciles.

\subsection{Resumen de la sección}

Para Xie Saining, LeCun es importante no solo por su posición histórica, sino porque unifica en una sola persona la personalidad, la gestión y el marco metodológico. La amplitud de JEPA, la cautela frente a la idea de los LLM como destino final, la gestión al estilo del velero y la capacidad de impulso propia de una «batería enorme» explican en conjunto por qué Xie Saining está dispuesto a apostar la siguiente etapa de su trabajo a acompañarlo.

